% Source: sources/returns/S6_04_MAPS_MANIFOLDS/sections/manifolds.tex:3980–4289
\paragraph{Theorem D.2.20 (Proved: six original-filled families, including the untwisted Engel case).}
Use the following configurations and section choices:
\[
\begin{array}{c|c|l|r}
 \mathrm{OS}&\mathrm{profile}&\text{fibres}&d\\\hline
 43&\mathrm{default}&III^*,3I_1&2\\
 43&II&III^*,II,I_1&2\\
 46&\mathrm{default}&I_3^*,3I_1&4\\
 46&II&I_3^*,II,I_1&4\\
 49&III&IV^*,III,I_1&6\\
 50&IV&I_1^*,IV,I_1&12.
\end{array}
\]
Take $P=G$, $Q=dG$, retain every original additive filling,
and put the simple linearization zero at the first marked
$I_1$. The other $I_1$ fillings are unchanged. At an
additional regular fibre perform any free log
\[
 m>0,\qquad v=a e_1+b e_2+c\delta+n\gamma,
 \qquad\gcd(m,a,b,c,n)=1.
\]
Set
\[
        g=\gcd(m,a,b,n),\qquad
        T_{m,g}=\mathbb Z/g\oplus(\mathbb Z/m)^2.
\]
For $n\ne0$, the answer is
\[
\begin{aligned}
 \pi_1(X)&=H_1(X)=\mathbb Z/|n|, &H_5(X)&=0,\\
 H_2(X)&=\mathbb Z^7\oplus T_{m,g},
 &H_3(X)&=\mathbb Z^{14}\oplus T_{m,g},\\
 H_4(X)&=\mathbb Z^7\oplus\mathbb Z/|n|.
\end{aligned}
\tag{D.O1}
\]
For $n=0$, it is
\[
\begin{aligned}
 \pi_1(X)&=H_1(X)=H_5(X)=\mathbb Z,\\
 H_2(X)&=\mathbb Z^8\oplus T_{m,g},
 &H_3(X)&=\mathbb Z^{14}\oplus T_{m,g},\\
 H_4(X)&=\mathbb Z^8.
\end{aligned}
\tag{D.O2}
\]
In all cases $H_0=H_6=\mathbb Z$ and $\chi(X)=2$.
The simply connected members are exactly $n=\pm1$, and
all of them have oriented diffeomorphism type
\[
 X\cong\#^7(S^2\times S^4)\#^7(S^3\times S^3)
                                    \#\widetilde\Sigma_0N_m.
 \tag{D.O3}
\]
In particular none is a six-sphere. For the Engel row,
the trivial regular log $(m,v)=(1,0)$ recovers the
untwisted K13 groups $(1,1,8,14,8,1,1)$ integrally,
with no torsion.
\begin{proof}
\emph{The original additive bundles.}
The section $Q$ is narrow everywhere. At every additive
fibre except $II$, $P$ meets a nonidentity component.
Proposition D.2.19 supplies the degree-one linearization.
The zero on $I_1$ gives the same primitive hexagonal
cusp as before. This proof retains the original additive
bundles, including their nonzero component degrees.

Let $E\to\Delta$ be one such additive elliptic-surface
neighbourhood, with $r$ components. It is simply connected
and retracts onto a connected union of $r$ rational spheres
with $H_2(E)=\mathbb Z^r$ and no other positive homology.
This applies also to the tangent pair of type $III$ and
the three spheres of type $IV$, which each meet at one
point. The component degree vector of $M=\mathcal O(P-O)$
is the primitive vector
\[
                   c=e_{\bar P}-e_0\in H^2(E;\mathbb Z).
\]
Because $Q$ is narrow, its local logarithm contracts it
to zero through translations. Rigidify the line along
$O$ as in D.2.16. The unit in the lift has a logarithm
since $H^1(E)=0$, and can be straightened. The actual
original filling is therefore, with its transported
boundary marking, the principal $T^2$ bundle over $E$
whose two Chern classes are $(c,0)$. Its two circle
directions are denoted $\delta,\gamma$.

The integral bundle spectral sequence has only base
rows zero and two. Its nonzero differentials are
\[
 H_2(E)\longrightarrow H_1(T^2),\quad x\longmapsto c(x)\delta,
 \qquad
 H_2(E)\otimes H_1(T^2)\longrightarrow H_2(T^2),\quad
 x\otimes y\longmapsto c(x)\delta\wedge y.
\]
Both images are primitive because $c$ is primitive.
There are no higher differentials, and all surviving
groups and extensions are free. Consequently
\[
 \pi_1(U)=\mathbb Z\gamma,\qquad
 (H_0,H_1,H_2,H_3,H_4)(U)
  =(\mathbb Z,\mathbb Z,\mathbb Z^{r-1},
                             \mathbb Z^{2r-1},\mathbb Z^r).
 \tag{D.O4}
\]
In particular $H_2(U)\to H_2(E)$ identifies its image
with $\ker c$.

\emph{The marked fibre kernel.}
Write the upper-left elliptic monodromy block as $B_i$
and its $\gamma$ column as $q_i$. Local narrowness of $Q$
means that
\[
                 q_i=(B_i-I)z_i,\qquad z_i\in\mathbb Z^2.
\]
The logarithm deformation identifies the elliptic
projection of a regular four-torus by
$e_1\mapsto e_1$, $e_2\mapsto e_2$, $\delta\mapsto0$,
$\gamma\mapsto z_i$. Thus its area covector, in the
usual ordered exterior basis, is
\[
           w_i=(1,0,(z_i)_2,0,-(z_i)_1,0).
 \tag{D.O5}
\]
The specialization of the elliptic fibre in $H_2(E)$
is its component-multiplicity vector. This vector is
primitive, lies in $\ker c$, and generates the image
of $H_2(F)\to H_2(U)$. The latter map is exactly
$w_i$ times that primitive vector, by (D.O4).
Therefore its kernel is the full integral $\ker w_i$.

The boundary monodromies in these six cases have two
unit nonzero Smith entries in $A_i-I$ and four in
$\bigwedge^2A_i-I$. Hence $H_1(\partial U)=\mathbb Z^3$
and $H_2(\partial U)=\mathbb Z^4$.
The two principal circle actions sweep the origin
disc, giving normal tori with orbit vectors
$\delta$ and $(-z_i,0,1)$. They form the full fixed
lattice of $A_i$. Thus these two sweeps split the
Wang projection integrally, and the remaining
boundary image equals the regular-fibre image just
computed. It follows that the local free defect is
\[
                H_2(U)/\operatorname{im}H_2(\partial U)
                         =\mathbb Z^{r-2}.
 \tag{D.O6}
\]
The full boundary kernel consists of the two actual
normal tori and the fibre kernel (D.O5). This is why
using only the fixed-loop normal tori would lose
relations in this case.

The exterior coinvariants of the global regular family
are free of rank one, with primitive detector
$\omega=e_1^*e_2^*-d\delta^*\gamma^*$.
If $s$ has $\omega$-value one, the image of $\ker w_i$
in this group is $h_i\mathbb Zs$, where
\[
                  h_i=\gcd(d,(z_i)_1,(z_i)_2).
\]
Indeed an integral basis of $\ker w_i$ has $\omega$
values $0,-(z_i)_2,0,(z_i)_1,-d$.
Here all local inputs are explicit:
\[
\begin{array}{c|c|r|r|r}
 \mathrm{OS/profile}&\text{non-narrow-}P\text{ fibre}
       &r&z_i&h_i\\\hline
 43/\mathrm{default},43/II&III^*&8&(3,1)&1\\
 46/\mathrm{default},46/II&I_3^*&8&(1,-2)&1\\
 49/III&IV^*&7&(6,2)&2\\
       &III&2&(9,3)&3\\
 50/IV&I_1^*&6&(9,-6)&3\\
      &IV&3&(8,-8)&4.
\end{array}
\tag{D.O7}
\]
In every row of the theorem the $h_i$ together have gcd
one. Thus the actual original fillings kill $s$.
The original $II$ and ordinary $I_1$ pieces have the
product models already proved and add no defect.
The total free local defect $\beta$ is six in the first
four configurations and five in the last two.

\emph{The unrestricted global matrices.}
Insert $A_0=I$ for the regular log. If the total number
of punctures is $N$, the graph differential has domain
$\mathbb Z^{4(N-1)}$, rank three and primitive image.
Write $B$ for all unchanged fixed-loop normal columns,
including the two at the cusp. The exact constant
certificates give an integral basis
\[
 K=[B\ j_0(e_1)\ j_0(e_2)\ j_0(\gamma)]
\]
of the graph kernel, with left inverse $L$ whose last
three coordinates on $j_0(x)$ are $(x_1,x_2,x_4)$.
The actual circle sweeps above fix their zero fibre-area
coordinates. After these normal columns and the relations
$h_is=0$ are eliminated, the regular-family contribution
is freely generated by those three sweeps.

The smooth log has lattice $\Lambda+\mathbb Z(v/m)$,
with primitive relation $m(v/m)=v$.
Its normal tori have graph components $j_0(mx)$;
the additional lattice generator has component $j_0v$.
The fibre components no longer matter, since $s=0$.
The resulting presentation is exactly
\[
  \begin{pmatrix}
    m&0&0&a\\0&m&0&b\\0&0&m&n
  \end{pmatrix}.
 \tag{D.O8}
\]
Its determinantal divisors are $g,mg,m^2g$,
so its cokernel is $T_{m,g}$ for all integers in the
statement. Notice that $c$ enters the geometric
freeness condition even though it drops out of (D.O8).

The original fillings kill all base meridians and
precisely the primitive $\ker\gamma^*$ in the fibre
fundamental group. Their complement over the punctured
base therefore has fundamental group $\mathbb Z\gamma$.
The new meridian relation is $m\xi-v=0$, and the base
relation kills $\xi$, leaving $n\gamma=0$. This proves
the asserted full fundamental group.

The degree-one Mayer--Vietoris map must now be counted
with the groups (D.O4), not with the rank-two groups of
the narrow-$P$ filling. Its source and target are free.
For the first and third rows of the theorem their ranks
are $(19,18)$; for the second and fourth, $(14,13)$;
for the last two, $(14,12)$.
Its cokernel is the abelian group already computed by
van Kampen. Thus for $n\ne0$ its kernel rank is
$\varepsilon=1$ in the first four rows and two in the
last two. For $n=0$ this rank increases by one.
These free kernels split off in the degree-two sequence.
Since $\beta+\varepsilon=7$, this proves exactly the
stated $H_2$ groups, including their torsion.
All local pieces except the primitive cusp have Euler
zero, while the cusp contributes two. Integral duality
then gives (D.O1)--(D.O2) in every degree. In particular
the Engel row at $(m,v)=(1,0)$ independently reproduces
K13, including its integral freeness.

\emph{Canonical class, nullsphere surgery and smooth type.}
Let $F_m$ denote the reduced new multiple fibre.
The canonical-form computation in D.3.3 applies to the
original pieces: before the regular log the canonical
line is $f^*\mathcal O(-1)$, and the primitive toroidal
cusp adds no divisor. A regular order-$m$ log adds
$(m-1)F_m$. Since $f^*(\mathrm{point})=mF_m$, one obtains
\[
                         K_X=\mathcal O_X(-F_m).
 \tag{D.O9}
\]
Here this formula also controls the topological class,
including the free classes supplied by the degree-one
Mayer--Vietoris kernel. The complement of $F_m$ has
$\pi_1=\mathbb Z\gamma$. Its primitive normal meridian
is $m\xi-v$, whose image is $-n\gamma$.
The residue sequence for the divisor therefore shows
that $[F_m]\in H^2(X;\mathbb Z)$ has exact order $|n|$
when $n\ne0$. For $n=0$ it is primitive of infinite
order: the quotient by its image embeds in the torsion-free
$H^2(X\setminus F_m)$, whose torsion vanishes by
$H_1(X\setminus F_m)=\mathbb Z$. For $n=\pm1$ it vanishes. Thus $c_1(X)=0$ in the
simply connected cases, and $X$ is spin; local canonical
triviality alone would not have sufficed for this step.

The angular $\delta$ action is free on every original
bundle and on the smooth log when $n=\pm1$. Its sole
fixed curve is the cusp rational curve with normal
weights $(1,-1)$. This curve belongs to the regular-fibre
specialization image and hence is a multiple of $s$.
By (D.O7), $s=0$ already in the exterior after the
original additive fillings. The fixed curve is therefore
null-homologous, and then null-homotopic in a simply
connected member.
Apply the oriented equivariant nullsphere surgery and
projectable spin lift of D.2.14. It gives
$Y=X\#(S^3\times S^3)$ as a free circle bundle over a
simply connected spin five-manifold with
\[
            H_2(B_5)=\mathbb Z^8\oplus(\mathbb Z/m)^2
\]
and primitive Euler class. Smale--Barden identifies the
oriented pair as
\[
 (\#^8(S^2\times S^3)\#N_m,(x,0,\ldots,0)).
\]
Repeated \cite[Theorem A and Theorem B(2)]{manGR25}
give
\[
 Y\cong\#^7(S^2\times S^4)\#^8(S^3\times S^3)
                                        \#\widetilde\Sigma_0N_m.
\]
Cancellation by the complete Zhubr invariants proves
(D.O3), including the specified twisted suspension
framing at even $m$.

The exact program \texttt{manifolds\_original\_additive.py}
retains all monodromies, fixed-loop bases, local covectors,
primitive fibre relations and full integral gluing maps.
It verifies 4,386 full degree-two Smith forms and 4,386
full degree-one Smith forms, including 582 cases with
$n=0$ in each degree. It separately checks the base-degree
map on bridge classes in both rank regimes. Formula
(D.O8) and the proved integral local maps supply the
unrestricted theorem. No priority claim is made.
\end{proof}
